% ============================================================================
%  Rozdział 4. Implementacja systemu
% ============================================================================
\chapter{Implementacja systemu}
\label{rozdz:implementacja}

W~poprzednich rozdziałach przedstawiono założenia, wymagania i~projekt architektury.
Ten rozdział opisuje, w~jaki sposób zostały one zrealizowane w~kodzie. Omówiono kolejno
strukturę projektu, wczytywanie konfiguracji, kompozycję zależności, sterowanie i~model
ruchu kuli, przeszkody, przeciwników i~punkty życia, monety i~warunek ukończenia poziomu,
pulę efektów cząsteczkowych, interfejs użytkownika w~podejściu MVVM, animacje
asynchroniczne, minimapę, pomiar czasu oraz walidację konfiguracji przed budowaniem. Fragmenty
kodu pochodzą bezpośrednio z~repozytorium projektu; pominięto w~nich dyrektywy
\kod{using}, deklaracje przestrzeni nazw oraz części niezwiązane z~omawianym zagadnieniem.

\section{Struktura projektu}

Kod gry znajduje się w~katalogu \sciezka{Assets/Code} i~liczy 256~plików C\# o~łącznej
objętości około 7~tys. wierszy. Zasoby Unity (sceny, prefabrykaty, materiały, animacje)
oraz pliki konfiguracyjne są przechowywane w~osobnych katalogach, a~kod narzędzi edytora
w~katalogu \sciezka{Assets/Editor}. Biblioteki
zewnętrzne znajdują się w~katalogach \kod{Plugins}, \kod{Appodeal}, \kod{Firebase}
i~\kod{Packages}. Najważniejsze elementy struktury przedstawiono na
rysunku~\ref{rys:struktura}.

\begin{figure}[!htb]
  \centering
  \begin{minipage}{0.92\textwidth}
  \dirtree{%
    .1 Assets/.
    .2 Code/\DTcomment{kod gry podzielony na 50 zestawów kompilacji}.
    .3 Configuration/\DTcomment{klasy konfiguracji i~danych, link.xml}.
    .3 Infrastructure/\DTcomment{instalatory, bindery i~fabryki Zenject}.
    .3 Main Scene/\DTcomment{kula, monety, przeciwnicy, przeszkody, interfejs}.
    .3 Main Menu/\DTcomment{piłka w~tle menu, okno statystyk}.
    .3 Loading/, Bootstrap/, All Scenes/\DTcomment{przepływ scen i~statystyki}.
    .3 Firebase Settings/, Ad Settings/\DTcomment{usługi zewnętrzne}.
    .2 Editor/\DTcomment{walidacja przed budowaniem, konfiguracja poziomów}.
    .2 Input/\DTcomment{akcje Input System i~odczyt pozycji wskaźnika}.
    .2 Resources/.
    .3 Configurations/\DTcomment{18 plików JSON z~parametrami rozgrywki}.
    .3 ProjectContext.prefab\DTcomment{kontekst projektu Zenject}.
    .2 Scenes/\DTcomment{Bootstrap, Main Menu, Loading, Level 1--3}.
    .2 Prefabs/\DTcomment{kula, monety, przeciwnicy, przeszkody, elementy HUD}.
    .2 Plugins/\DTcomment{Zenject, biblioteki natywne}.
  }
  \end{minipage}
  \caption{Struktura katalogów projektu; \zrodloWlasne}
  \label{rys:struktura}
\end{figure}

Kod podzielono na 50~zestawów kompilacji, a~dodatkowy zestaw \kod{Input} zawiera klasę
akcji wygenerowaną przez pakiet Input System. Zestawy deklarują łącznie 136~odwołań do
innych zestawów, zarówno projektu, jak i~bibliotek takich jak Zenject czy UniTask. 11~z~nich
nie odwołuje się do żadnego innego zestawu (są to głównie komponenty znacznikowe i~klasy
danych), a~najwięcej zależności (27) ma zestaw \kod{Main-Scene-Installer}, co jest zgodne
z~jego rolą korzenia kompozycji sceny poziomu. Podział zestawów na grupy przedstawia tabela~\ref{tab:asmdef}.

\begin{table}[!htb]
  \centering
  \caption{Grupy zestawów kompilacji w~projekcie; \zrodloWlasne}
  \label{tab:asmdef}
  \begin{tabularx}{\textwidth}{>{\raggedright\arraybackslash}p{3.2cm} c X}
    \toprule
    \textbf{Grupa} & \textbf{Liczba} & \textbf{Zestawy} \\
    \midrule
    Infrastruktura & 6 & Main-Scene-Installer, Main-Menu-Installer, Loading-Installers, Bootstrap-Installer, All-Scenes-Installer, Load-Level \\
    Konfiguracja & 1 & Configuration \\
    Kula i~sterowanie & 10 & Ball, Ball-Colliders, BallDestroyer, Ball-Marker, Bounce-Animation, UI-Displayable, User-Input, Ray Points, Line-Render, Main-Camera \\
    Obiekty planszy & 17 & Coin, Coin-Value, Enemy, Finish, Finish-Marker, Finish-Trigger, Ground, Wall, WallsMarker, Heal, Boost-Pad, Obstacles, Wind, Objects, StartPoints, Classes-Parameters, Dust-Effect \\
    Interfejs & 10 & UI-Coins, UI-Health-Bar, Information-Panel, Level-Timer, Minimap, Push-Button, Finish-Window, Finish-Window-Marker, UIAnimations, Buttons \\
    Sceny & 5 & All-Scenes, Bootstrap, Loading, Main-Menu, Main-Menu-Ball \\
    Usługi & 1 & Firebase-Settings \\
    \bottomrule
  \end{tabularx}
\end{table}

\section{Konfiguracja parametrów rozgrywki}

Wszystkie wartości liczbowe wpływające na rozgrywkę są przechowywane w~plikach JSON
w~katalogu \sciezka{Resources/Configurations}. Plik jest podzielony na sekcje, a~każdej
sekcji odpowiada para klas: klasa danych z~polami \kod{readonly} i~konstruktorem
przyjmującym wszystkie wartości oraz klasa konfiguracji dziedzicząca po generycznej klasie
\klasa{AbstractConfiguration<T>} (listing~\ref{lst:konfiguracja}). Klasa bazowa wczytuje
plik jako zasób tekstowy, wyodrębnia wskazaną sekcję i~deserializuje ją biblioteką
Newtonsoft.Json \cite{newtonsoft}, która dopasowuje nazwy pól JSON do parametrów
konstruktora.

\begin{lstlisting}[caption={Generyczna klasa wczytująca sekcję pliku konfiguracyjnego; \zrodloWlasne},label=lst:konfiguracja]
public abstract class AbstractConfiguration<T>
{
    private readonly string _resourcePath, _subsectionName;
    protected AbstractConfiguration(string jsonName, string subsectionName)
    {
        _resourcePath = "Configurations/" + jsonName;
        _subsectionName = subsectionName;
    }
    public T Load()
    {
        TextAsset config = Resources.Load<TextAsset>(_resourcePath);
        if (config == null)
            throw new FileNotFoundException($"Config resource not found: {_resourcePath}");
        JToken subsection = JObject.Parse(config.text)[_subsectionName];
        if (subsection == null)
            throw new InvalidDataException($"Config {_resourcePath} has no section '{_subsectionName}'");
        return subsection.ToObject<T>();
    }
}

public class BallMovementConfiguration : AbstractConfiguration<BallMovementData>
{
    public BallMovementConfiguration() : base("Ball Configuration", "Movement") { }
}
\end{lstlisting}

Klasy danych, np. \klasa{BallMovementData} z~polami \kod{MaxPushingPower}
i~\kod{PushPowerDecrease}, nie mają żadnej logiki i~nie mogą zostać zmienione po
utworzeniu, co wyklucza przypadkową modyfikację parametrów w~trakcie gry. Ponieważ
Newtonsoft.Json tworzy obiekty przez refleksję, kompilator IL2CPP mógłby uznać ich
konstruktory za nieużywane i~usunąć je z~wersji wynikowej. Zapobiega temu plik
\kod{link.xml} z~wpisem \kod{<assembly fullname="Configuration" preserve="all"/>},
który wyłącza z~usuwania cały zestaw konfiguracji. Komplet plików konfiguracyjnych
zamieszczono w~załączniku~B.

\section{Kompozycja zależności z wykorzystaniem Zenject}
\label{sec:kompozycja}

Kontekst projektu zawiera tylko jeden instalator, który rejestruje obiekt
\klasa{LoadLevel} przechowujący nazwę poziomu wybranego w~menu. Pozostałe powiązania są
rejestrowane w~kontekstach scen. Instalatory scen poziomów nie wywołują metod kontenera
bezpośrednio, lecz budują listę binderów, czyli obiektów klasy pochodnej
\klasa{Binder} (listing~\ref{lst:binder}). Każdy binder rejestruje jedną spójną grupę
powiązań, np. dane wiatru i~obiekt \klasa{LevelWind} losujący kierunek oraz siłę wiatru
na dany poziom.

\begin{lstlisting}[caption={Klasa bazowa binderów i~binder rejestrujący wiatr; \zrodloWlasne},label=lst:binder]
public abstract class Binder
{
    protected readonly DiContainer Container;
    protected Binder(DiContainer container) => Container = container;
    public abstract void Bind();
}

public class BindToWind : Binder
{
    private readonly WindData _windData;
    public BindToWind(DiContainer container, WindData windData)
        : base(container) => _windData = windData;
    public override void Bind()
    {
        Container.Bind<WindData>().FromInstance(_windData).AsSingle();
        Container.Bind<LevelWind>().FromInstance(new LevelWind(_windData)).AsSingle();
    }
}
\end{lstlisting}

Najciekawszy przypadek stanowi instalator kuli (listing~\ref{lst:ball-installer}).
Kula nie istnieje w~scenie, lecz jest tworzona z~prefabrykatu podczas instalacji
powiązań. Instalator działa w~dwóch etapach. Najpierw rejestruje zależności potrzebne
kuli (dane ruchu, animacji, zdrowia, wiatru, komponent wejścia i~wyzwalacz mety),
a~następnie tworzy kulę przez fabrykę korzystającą z~\kod{DiContainer.InstantiatePrefab},
dzięki czemu wszystkie jej komponenty otrzymują te zależności. W~drugim etapie binder
\klasa{BindFromBall} rejestruje w~kontenerze komponenty utworzonej kuli
(\klasa{BallMarker}, \klasa{HealthBarModel}, \klasa{BallSpeedSender} i~inne), aby mogły
z~nich korzystać kamera, minimapa, linia celownicza i~elementy interfejsu.

\begin{lstlisting}[caption={Dwuetapowa instalacja powiązań kuli w~klasie BallInstaller; \zrodloWlasne},label=lst:ball-installer]
public override void InstallBindings()
{
    CreateToBallBinders();
    BindToBall();         // zależności potrzebne kuli
    CreateFromBallBinders();
    BindFromBall();       // komponenty utworzonej kuli dla pozostałych modułów
}
private void CreateFromBallBinders()
{
    BallFactory ballFactory = new
    (
        Container,
        new BallPrefab(_ballMarker.gameObject),
        _ballStartPoint.position,
        Quaternion.identity
    );
    _fromBallBinders.Add(new BindFromBall(Container, ballFactory.InstantiateBall()));
}
\end{lstlisting}

Komponenty \klasa{MonoBehaviour} otrzymują zależności przez metodę oznaczoną atrybutem
\kod{[Inject]}, nazwaną w~projekcie \kod{Constructor}. Typowy komponent rozgrywki
(listing~\ref{lst:ball-movement}) zapisuje otrzymane obiekty w~polach, w~metodzie
\kod{Start} tworzy obiekt logiki, w~\kod{FixedUpdate} wywołuje jego metodę, a~przy
dezaktywacji zwalnia jego subskrypcje. Sam komponent nie zawiera żadnej reguły gry.

\begin{lstlisting}[caption={Komponent BallMovement jako adapter pomiędzy silnikiem a~logiką ruchu; \zrodloWlasne},label=lst:ball-movement]
[RequireComponent(typeof(Rigidbody), typeof(SlowDownBall), typeof(BounceBall))]
public class BallMovement : MonoBehaviour, IBallMovementContainer
{
    public AbstractBallMovementLogic BallMovementLogic { get; private set; }
    private BallInput _ballInput;
    private BallMovementData _ballMovementData;
    [Inject]
    private void Constructor(BallMovementData ballMovementData, BallInput ballInput)
    {
        _ballMovementData = ballMovementData;
        _ballInput        = ballInput;
    }
    private void Start()
    {
        BallMovementLogic = new BallMovementLogic
        (
            GetComponent<Rigidbody>(), transform, _ballInput,
            GetComponent<BounceBall>(), GetComponent<SlowDownBall>(),
            _ballMovementData.MaxPushingPower, _ballMovementData.PushPowerDecrease
        );
    }
    private void FixedUpdate() => BallMovementLogic.PushBall();
    private void OnDisable() => BallMovementLogic.Dispose();
}
\end{lstlisting}

Prefabrykaty przekazywane fabrykom są opakowane w~klasy takie jak \klasa{CoinPrefab},
\klasa{CactusPrefab} czy \klasa{Walls}, dziedziczące po wspólnej klasie bazowej
\klasa{AbstractGameObjectParametr}. Dzięki temu kontener rozróżnia kilka obiektów typu
\klasa{GameObject} po typie opakowania, bez identyfikatorów tekstowych.

\section{Sterowanie kulą}
\label{sec:sterowanie}

Akcja wejściowa \kod{Move} pliku \kod{Player Input.inputactions} jest typu przycisk
i~ma dwa powiązania: lewy przycisk myszy oraz naciśnięcie pierwszego dotyku na ekranie.
Wciśnięcie wywołuje zdarzenie \kod{performed}, a~zwolnienie \kod{canceled}. Pozycję
wskaźnika odczytuje statyczna klasa \klasa{PlayerPointer} z~uogólnionego wskaźnika
\klasa{Pointer}, obsługującego mysz. Wariant odczytujący pozycję pierwszego dotyku
z~urządzenia \klasa{Touchscreen} jest wybierany dyrektywą kompilacji warunkowej.

Po wciśnięciu klasa \klasa{BallInputLogic} rzutuje promień z~kamery przez punkt ekranu
i~sprawdza, czy pierwszy trafiony zderzacz należy do kuli. Wywołanie \kod{Physics.Raycast}
otrzymuje parametr \kod{QueryTriggerInteraction.Ignore}, więc pole przyspieszające lub
apteczka, w~których stoi kula, nie przesłaniają jej. Tylko gest rozpoczęty na kuli
generuje zdarzenia wciśnięcia i~zwolnienia przekazywane dalej. Zdarzenie wciśnięcia
aktywuje linię celowniczą (komponent \klasa{LineRenderer} rysowany w~każdym kroku fizyki
od kuli do punktu podłoża pod kursorem) oraz zapamiętuje punkt początkowy gestu
(rysunek~\ref{rys:celowanie}).

\begin{figure}[!htb]
  \centering
  \zrzut[0.66\textwidth]{poziom1_celowanie}
  \caption{Celowanie: linia od kuli do punktu podłoża pod kursorem; po zwolnieniu kula
    porusza się w~kierunku przeciwnym; \zrodloWlasne}
  \label{rys:celowanie}
\end{figure}

Moc i~kierunek strzału wyznacza klasa \klasa{PushPowerConverter}
(listing~\ref{lst:push-power}). Moc zależy od odległości w~pikselach pomiędzy punktem
wciśnięcia $\mathbf{s}_p$ i~zwolnienia $\mathbf{s}_r$ na ekranie, a~kierunek od położenia
punktu zwolnienia na płaszczyźnie podłoża $\mathbf{g}_r$ względem rzutu środka kuli
$\mathbf{b}$ na tę płaszczyznę:
\begin{equation}
  P = \min\left(\frac{\lVert \mathbf{s}_r - \mathbf{s}_p \rVert}{k},\; P_{\max}\right), \qquad
  \mathbf{d} = -\frac{\mathbf{g}_r - \mathbf{b}}{\lVert \mathbf{g}_r - \mathbf{b} \rVert},
  \label{eq:strzal}
\end{equation}
gdzie $k = 20$~px jest dzielnikiem mocy, a~$P_{\max} = 25$ mocą maksymalną. Maksymalną
moc uzyskuje się więc przy przeciągnięciu o~500~pikseli. Punkt $\mathbf{g}_r$ jest
wyznaczany jako przecięcie promienia kamery z~płaszczyzną $y = 0$ metodą
\kod{Plane.Raycast}, co jest tańsze od rzutowania promienia na zderzacze sceny.

\begin{lstlisting}[caption={Wyznaczanie mocy i~kierunku strzału; \zrodloWlasne},label=lst:push-power]
public (float distance, Vector3 direction) ConvertToPushPower()
{
    float distance = Vector3.Distance(ScreenStartPoint, ScreenFinishPoint) / _pushPowerDecrease;
    Vector3 normalizedDirection = -(GroundFinishPoint - GroundStartPoint).normalized;
    if (distance > _maxPushingPower) distance = _maxPushingPower;
    return (distance, normalizedDirection);
}
\end{lstlisting}

\section{Model ruchu, tarcie i odbicia}

Stan ruchu kuli opisują dwie wielkości: jednostkowy wektor kierunku $\mathbf{d}$
i~moc $P$. W~każdym kroku fizyki logika ruchu nadpisuje prędkość ciała sztywnego:
\begin{equation}
  \mathbf{v} = P\,\mathbf{d} + \mathbf{w},
  \label{eq:predkosc}
\end{equation}
gdzie $\mathbf{w}$ jest prędkością dryfu wywołanego wiatrem. Ponieważ kierunek leży
zawsze w~płaszczyźnie poziomej, kula nie odrywa się od podłoża, a~moc $P$ jest
jednocześnie jej prędkością w~metrach na sekundę. Klasa bazowa logiki ruchu udostępnia
trzy operacje zmiany stanu (listing~\ref{lst:ruch}): \kod{Launch} dodaje moc z~ograniczeniem
do zadanej wartości i~jest używana przez pola przyspieszające, odbijacze oraz przycisk
\emph{Push}, \kod{Knock} gwarantuje jedynie minimalną moc i~jest używana przez przeszkody
ruchome, a~prywatna metoda \kod{BounceBall} obsługuje odbicia. Każda zmiana mocy jest
publikowana w~strumieniu \kod{PushingPowerChanged}, z~którego korzysta panel prędkości.

\begin{lstlisting}[caption={Operacje zmiany stanu ruchu w~klasie AbstractBallMovementLogic; \zrodloWlasne},label=lst:ruch]
// Drift is added on top of the push and is not slowed by friction, so it also moves a ball at rest.
public void PushBall() => _rigidbody.velocity = Direction * PushingPower + Drift;
public void Launch(Vector3 direction, float power, float maxPushingPower)
{
    Direction    = direction;
    PushingPower = Mathf.Min(PushingPower + power, maxPushingPower);
    PushingPowerChanged.OnNext(PushingPower);
    BallPushed?.Invoke(Direction);
}
// Unlike Launch it adds no speed: a fast ball keeps its speed, a slow one gets at least the given power.
public void Knock(Vector3 direction, float power)
{
    Direction    = direction;
    PushingPower = Mathf.Max(PushingPower, power);
    PushingPowerChanged.OnNext(PushingPower);
    BallPushed?.Invoke(Direction);
}
private void BounceBall(Vector3 bounceDirection, float lostPower)
{
    Direction = bounceDirection;
    PushingPower *= lostPower;
    PushingPowerChanged.OnNext(PushingPower);
    BallPushed?.Invoke(Direction);
}
\end{lstlisting}

Tarcie realizuje komponent \klasa{Friction} podłoża. W~metodzie \kod{OnCollisionStay},
wywoływanej w~każdym kroku fizyki, w~którym kula styka się z~podłożem, zmniejsza on moc
o~stałą wartość:
\begin{equation}
  P_{n+1} = \max\left(P_n - f\,\Delta t,\; 0\right),
  \label{eq:tarcie}
\end{equation}
gdzie $f = 4$ jest siłą tarcia z~pliku konfiguracyjnego, a~$\Delta t = 0{,}02$~s.
Jest to ruch jednostajnie opóźniony, więc kula wystrzelona z~mocą maksymalną zatrzymuje
się po czasie $P_{\max}/f = 6{,}25$~s, przebywając drogę $P_{\max}^2/(2f) \approx 78$~m,
o~ile nie odbije się od przeszkody. Wartość ta jest porównywalna z~wymiarami poziomów
(tabela~\ref{tab:poziomy}), co oznacza, że najsilniejszy strzał pozwala przejechać przez
większą część planszy.

Odbicia obsługuje klasa \klasa{Bounce}, używana przez ściany i~przeciwników. Przy kolizji
z~kulą pobiera ona wektor normalny pierwszego punktu styku $\mathbf{n}$ i~przekazuje go
do logiki odbicia kuli (listing~\ref{lst:odbicie}), która wyznacza nowy kierunek zgodnie
z~prawem odbicia i~zmniejsza moc o~współczynnik $\lambda$ zależny od obiektu:
\begin{equation}
  \mathbf{d}' = \mathbf{d} - 2\,(\mathbf{d}\cdot\mathbf{n})\,\mathbf{n}, \qquad
  P' = \lambda P, \qquad \lambda_{\text{ściana}} = 0{,}7,\;\; \lambda_{\text{przeciwnik}} = 0{,}2.
  \label{eq:odbicie}
\end{equation}
Niska wartość współczynnika dla przeciwników sprawia, że zderzenie z~nimi prawie zatrzymuje
kulę, co stanowi dodatkową karę obok utraty zdrowia.

\begin{lstlisting}[caption={Wyznaczanie kierunku po odbiciu w~klasie BounceBallLogic; \zrodloWlasne},label=lst:odbicie]
public void GetBounceDirection(GameObject collidedObject, Vector3 normal, float lostPower)
{
    GetBounceDirectionAction?.Invoke
    (
        Vector3.Reflect(_ballMovement.BallMovementLogic.Direction, normal),
        lostPower
    );
    BallBounced?.Invoke(collidedObject);
}
\end{lstlisting}

Obrót siatki kuli jest liczony niezależnie od nadawanej prędkości. Klasa
\klasa{BallRotationLogic} w~każdym kroku fizyki wyznacza rzeczywiste przemieszczenie ciała
sztywnego i~obraca siatkę wokół osi $\mathbf{y}\times\hat{\mathbf{v}}$. Pierwsza wersja
obracała kulę na podstawie zadanej prędkości, przez co kula dociśnięta wiatrem do ściany
kręciła się w~miejscu; obecne rozwiązanie usunęło ten efekt.

\section{Przeszkody dynamiczne, pola specjalne i wiatr}

Pole przyspieszające ma zderzacz w~trybie wyzwalacza. Po wjechaniu kuli wywołuje operację
\kod{Launch} z~kierunkiem zgodnym z~osią $z$ pola rzutowaną na płaszczyznę poziomą,
z~dodatkiem mocy 20 i~ograniczeniem do 35. Kula jest więc wyrzucana wzdłuż strzałek
niezależnie od strony, z~której wjechała. Odbijacz działa na tej samej zasadzie, ale
kierunek wyznacza wektor od środka odbijacza do kuli. Po trafieniu korpus odbijacza
zostaje skokowo powiększony o~25\% i~w~ciągu 0{,}15~s liniowo wraca do pierwotnego
rozmiaru, co stanowi wizualne potwierdzenie zderzenia.

Obracające się belki i~ruchome ściany są ciałami kinematycznymi. Belka obraca się ze
stałą prędkością kątową 90°/s, a~kierunek obrotu jest losowany przy starcie poziomu.
Ruchoma ściana jest przesuwana w~każdym kroku fizyki metodą \kod{MovePosition} wzdłuż
własnej osi $x$ ruchem harmonicznym:
\begin{equation}
  \mathbf{x}(t) = \mathbf{c} + \mathbf{a}\,A \sin\!\left(\frac{v}{A}\,t + \varphi\right),
  \label{eq:sciana}
\end{equation}
gdzie $\mathbf{c}$ to położenie początkowe, $\mathbf{a}$ jednostkowy wektor osi, $A = 2{,}5$
amplituda, a~$\varphi$ losowa faza. Pochodna tego wyrażenia ma wartość maksymalną równą
$v$, więc parametr konfiguracyjny \kod{MovingWallSpeed} jest maksymalną prędkością
ściany, a~okres ruchu wynosi $2\pi A/v \approx 6{,}3$~s. Faza jest losowana w~metodzie
\kod{Awake}, dzięki czemu ściany na poziomie nie poruszają się synchronicznie.

Obie przeszkody mają komponent \klasa{ObstacleKnocker}, który przy kontakcie wywołuje
operację \kod{Knock} z~kierunkiem od punktu styku do środka kuli i~mocą co najmniej 10.
Bez niego kula mogła zostać zakleszczona pomiędzy ruchomą przeszkodą a~ścianą lub
przesuwana razem z~przeszkodą.

Wiatr jest losowany raz na poziom przez obiekt \klasa{LevelWind}: kierunek z~rozkładu
jednostajnego kąta w~zakresie $[0^\circ, 360^\circ)$, a~siła z~przedziału $[0{,}3;\,0{,}8]$.
Komponent \klasa{BallWind} w~każdym kroku fizyki ustawia prędkość dryfu $\mathbf{w}$
w~równaniu~(\ref{eq:predkosc}). Ponieważ dryf nie podlega tarciu, wiatr powoli przesuwa
również nieruchomą kulę. Kierunek wiatru pokazuje strzałka w~interfejsie, obracana
o~kąt pomiędzy kierunkiem wiatru a~rzutem osi kamery na płaszczyznę poziomą, dzięki czemu
wskazuje kierunek zgodny z~tym, co widzi gracz. Jej rozmiar rośnie liniowo wraz z~siłą
wiatru. Przycisk \emph{Push} jest aktywny tylko wtedy, gdy moc kuli wynosi zero, i~nadaje
jej moc 8 w~losowym kierunku.

\section{System przeciwników i punktów życia}

Punkty startowe przeciwników są obiektami potomnymi znacznika \klasa{AllEnemiesMarker}.
Każdy z~nich ma komponent \klasa{EnemyStartPoint} z~polem określającym typ przeciwnika.
Fabryka \klasa{EnemiesFactory} rozdziela punkty na dwie listy i~przekazuje je fabrykom
kaktusów i~grzybów. Fabryka kaktusów przekazuje do kontenera punkt startowy jako
dodatkowy argument wywołania \kod{InstantiatePrefabForComponent}, ponieważ dwa obiekty
potomne punktu startowego wyznaczają końce trasy patrolu. Logika ruchu kaktusa przesuwa go
metodą \kod{Vector3.MoveTowards} z~prędkością 3~m/s pomiędzy tymi punktami i~zawraca po
dotarciu do celu, a~osobna klasa obraca model w~kierunku bieżącego celu. Grzyby pozostają
nieruchome i~otrzymują losowy obrót wokół osi pionowej.

Prefabrykaty obu przeciwników mają dwa zderzacze potomne. \emph{Damage Collider} zadaje
obrażenia obiektowi, który z~nim kolidował, jeśli obiekt ten implementuje interfejs
\klasa{IGetDamage}. \emph{Death Collider} po zderzeniu z~kulą uruchamia animację śmierci,
zatrzymuje ruch kaktusa i~po 1,7~s dezaktywuje przeciwnika. Oba zderzacze odbijają kulę
ze współczynnikiem $\lambda_{\text{przeciwnik}}$. Interfejs \klasa{IGetDamage} implementuje
komponent \klasa{BallGetDamage}, który publikuje wartość obrażeń w~strumieniu R3. Strumień
ten subskrybuje model zdrowia \klasa{HealthBarModel}, delegujący reguły do klasy
\klasa{HealthStorage} (listing~\ref{lst:zdrowie}).

\begin{lstlisting}[caption={Przechowywanie punktów życia w~klasie HealthStorage; \zrodloWlasne},label=lst:zdrowie]
public class HealthStorage
{
    private const int _minHealth = 0;
    public Action Die;
    public readonly int DefaultHealth;
    public int Health { get; private set; }
    public HealthStorage(int health)
    {
        Health        = health;
        DefaultHealth = Health;
    }
    public void DecreaseHealth(int health)
    {
        Health -= health;
        if (Health < _minHealth)
        {
            Health = _minHealth;
            Die?.Invoke();
        }
    }
    public bool IncreaseHealth(int health)
    {
        if (Health >= DefaultHealth) return false;
        Health = Math.Min(Health + health, DefaultHealth);
        return true;
    }
}
\end{lstlisting}

Metoda \kod{IncreaseHealth} zwraca informację, czy leczenie zostało wykonane. Korzysta
z~niej apteczka: jeśli kula ma pełne zdrowie, apteczka pozostaje na planszy, w~przeciwnym
razie wyłącza swój zderzacz i~znika z~tą samą animacją co moneta. Zdarzenie \kod{Die}
modelu zdrowia jest przekazywane dalej jako strumień \kod{Subject<Unit>}; subskrybują go
komponent \klasa{BallDestroyer}, który dezaktywuje kulę i~publikuje zdarzenie zniszczenia
kuli, oraz licznik czasu, który zatrzymuje pomiar. Zdarzenie zniszczenia kuli powoduje
wyświetlenie okna końca poziomu (rysunek~\ref{rys:obrazenia}).

\begin{figure}[!htb]
  \centering
  \begin{minipage}{0.49\textwidth}
    \centering
    \zrzut[\textwidth]{poziom2_obrazenia}\\[2pt]
    {\small a)}
  \end{minipage}\hfill
  \begin{minipage}{0.49\textwidth}
    \centering
    \zrzut[\textwidth]{poziom2_przegrana}\\[2pt]
    {\small b)}
  \end{minipage}
  \caption{Drugi poziom: a) stan interfejsu po dwóch kolizjach z~kaktusem, b) okno końca
    poziomu po utracie zdrowia; \zrodloWlasne}
  \label{rys:obrazenia}
\end{figure}

\section{Monety i warunek ukończenia poziomu}
\label{sec:monety}

Monety nie są umieszczane w~scenie bezpośrednio. Projektant poziomu rozmieszcza puste
punkty startowe jako obiekty potomne znacznika \klasa{AllCoinsMarker}, a~fabryka monet
tworzy dla każdego z~nich monetę przez kontener (listing~\ref{lst:monety}). Fabryka
przechowuje monety na stosie, z~którego zdejmuje je przy rozmieszczaniu na planszy,
a~każda zebrana moneta wraca na stos po zakończeniu animacji znikania. Gdy liczba monet
na stosie ponownie zrówna się z~liczbą punktów startowych, wszystkie monety zostały
zebrane i~fabryka publikuje zdarzenie \kod{LastCoinCollected}. Rozwiązanie to nie wymaga
osobnego licznika ani przeszukiwania sceny.

\begin{lstlisting}[caption={Fragment fabryki monet; \zrodloWlasne},label=lst:monety]
private void FillCoins(Stack<CoinTransform> coins)
{
    foreach (CoinTransform _ in coins)
    {
        GameObject coin = InstantiateCoin();   // DiContainer.InstantiatePrefab
        _coins.Push(coin);
        SubscribeOnDestroying(coin.GetComponent<CoinDestroyer>());
    }
}
private void SubscribeOnDestroying(CoinDestroyer coinDestroyer)
{
    _coinsDestroyed.Push(coinDestroyer.CoinDestroyed.Subscribe(PushCoin));
}
private void PushCoin(GameObject coin)
{
    _coins.Push(coin);
    if (_coins.Count == _maxCoinsNumber) LastCoinCollected.OnNext(Unit.Default);
}
public void Dispose()
{
    while (!_coinsDestroyed.IsEmpty()) _coinsDestroyed.Pop().Dispose();
}
\end{lstlisting}

Pojedyncza moneta składa się z~kilku niezależnych komponentów: \klasa{CoinRotation}
obraca ją o~2° na krok fizyki, \klasa{CoinMovement} unosi ją i~opuszcza o~0,2~m,
\klasa{CoinValue} przechowuje wartość z~konfiguracji, a~\klasa{CoinTrigger} wykrywa
wjazd kuli. Wyzwalacz publikuje wartość monety w~strumieniu \kod{UpdateStorage} i~blokuje
ponowne zebranie. Fabryka rejestruje w~kontenerze zbiór strumieni wszystkich monet
(\klasa{CoinTriggersSubjects}), który subskrybują licznik monet w~interfejsie i~moduł
statystyk. Pełny łańcuch zdarzeń od zebrania monety do okna końca poziomu przedstawiono
na rysunku~\ref{rys:lancuch}, a~wszystkie strumienie reaktywne użyte w~grze zestawiono
w~tabeli~\ref{tab:strumienie}.

\begin{figure}[!htb]
  \centering
  \begin{tikzpicture}[diag]
    \node[modul, minimum width=2.45cm] (trig) at (0,0) {CoinTrigger\\ \kod{OnTriggerEnter}};
    \node[akcent, minimum width=2.45cm] (destr) at (3.75,0) {CoinDestroyerLogic\\ animacja (UniTask)};
    \node[akcent, minimum width=2.45cm] (fact) at (7.5,0) {CoinsFactory\\ \kod{PushCoin}};
    \node[akcent, minimum width=2.45cm] (fin) at (11.25,0) {FinishEnablerLogic\\ meta aktywna};
    \node[modul, minimum width=2.45cm] (ui) at (0,-1.7) {CoinsModel\\ licznik monet};
    \node[modul, minimum width=2.45cm] (stat) at (3.75,-1.7) {LevelStatisticsRecorder\\ maks. monet};
    \node[akcent, minimum width=2.45cm] (timer) at (7.5,-1.7) {LevelTimer\\ zapis wyniku};
    \node[modul, minimum width=2.45cm] (ftrig) at (11.25,-1.7) {FinishTrigger\\ wjazd na metę};
    \node[modul, minimum width=2.45cm] (window) at (3.75,-3.3) {FinishWindowShower\\ okno końca};
    \node[modul, minimum width=2.45cm] (destroy) at (7.5,-3.3) {BallDestroyer\\ kula nieaktywna};
    \draw[strz] (trig) -- node[above, font=\scriptsize]{Action} (destr);
    \draw[strz] (destr) -- node[above, font=\scriptsize]{Subject} (fact);
    \draw[strz] (fact) -- node[above, font=\scriptsize]{Subject} (fin);
    \draw[strz] (trig) -- node[left, font=\scriptsize]{Subject<int>} (ui);
    \draw[strz] (trig.south east) -- (stat.north west);
    \draw[strzp] (fin) -- node[right, font=\scriptsize]{wjazd kuli} (ftrig);
    \draw[strz] (ftrig) -- node[above, font=\scriptsize]{1} (timer);
    \draw[strz] (ftrig.south) |- node[pos=0.75, below, font=\scriptsize]{2} (destroy.east);
    \draw[strz] (destroy) -- (window);
  \end{tikzpicture}
  \caption{Łańcuch zdarzeń od zebrania monety do wyświetlenia okna końca poziomu; liczby
    oznaczają kolejność powiadamiania subskrybentów; \zrodloWlasne}
  \label{rys:lancuch}
\end{figure}

\begin{table}[!htb]
  \centering
  \caption{Strumienie reaktywne R3 wykorzystywane w~grze; \zrodloWlasne}
  \label{tab:strumienie}
  \footnotesize
  \begin{tabularx}{\textwidth}{>{\raggedright\arraybackslash}p{6.15cm} >{\raggedright\arraybackslash}p{5.25cm} X}
    \toprule
    \textbf{Strumień} & \textbf{Typ} & \textbf{Odbiorcy} \\
    \midrule
    \texttt{CoinTrigger.UpdateStorage} & \texttt{Subject<int>} & licznik monet, statystyki \\
    \texttt{CoinDestroyer.CoinDestroyed} & \texttt{Subject<GameObject>} & fabryka monet \\
    \texttt{CoinsFactory.LastCoinCollected} & \texttt{Subject<Unit>} & aktywacja mety \\
    \texttt{FinishTrigger.BallCollidedFinished} & \texttt{Subject<Unit>} & licznik czasu, kula \\
    \texttt{BallGetDamage.SendDamage} & \texttt{Subject<int>} & model zdrowia \\
    \texttt{HealthBarModel.Die} & \texttt{Subject<Unit>} & kula, licznik czasu \\
    \texttt{BallDestroyer.BallDestroyed} & \texttt{Subject<Unit>} & okno końca poziomu \\
    \texttt{PushingPowerChanged} & \texttt{Subject<float>} & nadajnik prędkości \\
    \texttt{BallSpeedSender.Speed} & \texttt{ReactiveProperty<float>} & panel prędkości \\
    \texttt{MeshRotation} & \texttt{ReactiveProperty<float>} & panel kąta \\
    \texttt{LevelTimerLogic.Result} & \texttt{ReactiveProperty<LevelResult>} & wynik w~oknie końca \\
    \bottomrule
  \end{tabularx}
\end{table}

Kolejność powiadomień ma tu znaczenie. Licznik czasu subskrybuje zdarzenie dotarcia do
mety już w~metodzie wstrzykiwania, a~komponent \klasa{BallDestroyer} dopiero
w~\kod{Start}. Ponieważ \klasa{Subject} powiadamia obserwatorów w~kolejności subskrypcji,
wynik poziomu jest zapisany, zanim kula zostanie dezaktywowana i~zanim okno końca poziomu
odczyta czas przejścia.

\section{Efekty cząsteczkowe i pula obiektów}

Każde odbicie kuli od ściany wyświetla w~punkcie styku krótki efekt pyłu oparty na
systemie cząsteczek. Tworzenie nowego obiektu efektu przy każdym odbiciu i~niszczenie go
po zakończeniu animacji powodowałoby ciągłe alokacje pamięci, dlatego zastosowano pulę
obiektów \cite{nystrom-pool} (listing~\ref{lst:pula}). Fabryka efektów przy starcie poziomu tworzy siedem
nieaktywnych egzemplarzy i~subskrybuje zdarzenie odbicia każdej ściany. Przy odbiciu
zdejmuje efekt ze stosu, aktywuje go i~ustawia w~punkcie styku zgodnie z~wektorem
normalnym. Po zakończeniu emisji Unity wywołuje metodę \kod{OnParticleSystemStopped}
komponentu efektu, który zgłasza to fabryce i~dezaktywuje się, a~fabryka odkłada go
z~powrotem na stos.

\begin{lstlisting}[caption={Pula efektów cząsteczkowych w~klasie DustEffectsFactory; \zrodloWlasne},label=lst:pula]
private const int _maxNumberOfDustEffects = 7;
private void Pop(Vector3 position, Vector3 normal)
{
    DustEffect dustEffect = _dustEffects.Pop();
    dustEffect.gameObject.SetActive(true);
    dustEffect.transform.position = position;
    dustEffect.transform.rotation = Quaternion.LookRotation(normal);
}
private DustEffect InstantiateDustEffect()
{
    DustEffect dustEffect = _diContainer.InstantiatePrefabForComponent<DustEffect>
        (_coinPrefab, Vector3.zero, Quaternion.identity, null);
    dustEffect.DustEffectStopped += PushDustEffect;
    dustEffect.gameObject.SetActive(false);
    return dustEffect;
}
private void FillDustEffectsStack()
{
    for (int i = 0; i < _maxNumberOfDustEffects; ++i) _dustEffects.Push(InstantiateDustEffect());
}
\end{lstlisting}

W~ten sposób w~trakcie rozgrywki nie powstaje ani nie jest niszczony żaden obiekt efektu,
a~liczba jednocześnie widocznych efektów jest ograniczona rozmiarem puli. Podobną rolę
pełni stos monet w~fabryce monet: monety są tworzone raz, przy wczytaniu poziomu,
a~zebrana moneta jest jedynie dezaktywowana.

\section{Interfejs użytkownika w architekturze MVVM}

Każdy element interfejsu, który wyświetla zmieniające się dane, jest zbudowany z~trzech
klas: modelu, modelu widoku i~widoku. Obiekty te tworzy komponent \klasa{MonoBehaviour}
dołączony do elementu interfejsu, który otrzymuje zależności przez wstrzykiwanie
i~zwalnia obiekty w~odwrotnej kolejności w~metodzie \kod{OnDestroy}. Dla panelu prędkości
i~kąta obrotu kuli przygotowano generyczne klasy bazowe (listing~\ref{lst:mvvm}): model
przepisuje wartość ze źródła, model widoku formatuje ją z~etykietą i~zaokrągleniem do
jednego miejsca po przecinku, a~widok ustawia tekst komponentu TextMeshPro. Konkretne
panele różnią się jedynie etykietą i~źródłem danych.

\begin{lstlisting}[caption={Generyczny model widoku i~widok panelu informacyjnego; \zrodloWlasne},label=lst:mvvm]
public abstract class AbstractViewModel : IDisposable
{
    public readonly ReactiveProperty<string> Text;
    private readonly string _standardTextPart;
    private readonly IDisposable _changeText;
    protected AbstractViewModel(ReactiveProperty<float> data, string standardTextPart)
    {
        _standardTextPart = standardTextPart;
        Text = new("");
        _changeText = data.Subscribe(ChangeText);
    }
    public void Dispose() => _changeText.Dispose();
    private void ChangeText(float data)
    {
        if (Text.Value == data.ToString()) return;
        Text.Value = _standardTextPart + " " + (float) Math.Round(value: data, digits: 1);
    }
}
public abstract class AbstractView : IDisposable
{
    private readonly TMP_Text _text;
    private readonly IDisposable _updateView;
    protected AbstractView(TMP_Text text, ReactiveProperty<string> textChanger)
    {
        _text = text;
        _updateView = textChanger.Subscribe(UpdateText);
    }
    private void UpdateText(string text) => _text.text = text;
    public void Dispose() => _updateView.Dispose();
}
\end{lstlisting}

Licznik monet ilustruje logikę prezentacji, której nie da się umieścić w~samym widoku.
Model subskrybuje strumienie wszystkich monet i~sumuje ich wartości. Model widoku nie
ustawia od razu nowej sumy, lecz zwiększa wyświetlaną liczbę o~jeden co 0,05~s
(listing~\ref{lst:licznik}), publikując przy każdym kroku zdarzenie, na które widok
reaguje pulsowaniem rozmiaru czcionki. Pasek zdrowia ma dwie pary model widoku--widok
korzystające z~tego samego modelu: dla suwaka i~dla liczby punktów życia.

\begin{lstlisting}[caption={Stopniowa zmiana licznika monet w~modelu widoku; \zrodloWlasne},label=lst:licznik]
public async UniTaskVoid Change(int target)
{
    int current = int.Parse(_coinsText.Value);
    if (current >= target) return;
    int k = current < target ? 1 : -1;
    while (current != target)
    {
        current += k;
        _coinsText.Value = current.ToString();
        PulsateText.OnNext(Unit.Default);
        await UniTask.WaitForSeconds(_changeSpeed);
    }
    FinishedTextPulsation.OnNext(Unit.Default);
}
\end{lstlisting}

\section{Animacje realizowane asynchronicznie}

Animacja odbicia kuli łączy trzy efekty: zmianę koloru na czerwony oraz, przy zderzeniu ze
ścianą, zgniecenie i~rozciągnięcie kuli. Każdy efekt implementuje interfejs
\klasa{IBallAnimation} z~metodą wykonującą jeden krok animacji i~metodą przywracającą stan
początkowy. Klasa \klasa{AnimationPlayer} wywołuje krok animacji w~każdym kroku fizyki
przez zadany czas, oczekując na \kod{UniTask.WaitForFixedUpdate}. Listing~\ref{lst:animacja}
pokazuje, jak z~tych elementów złożono sekwencję: zmiana koloru trwa 0,2~s i~przebiega
równolegle do dwóch następujących po sobie faz zgniatania i~rozciągania, trwających po
0,1~s. Szybkość deformacji jest proporcjonalna do prędkości kuli.

\begin{lstlisting}[caption={Sekwencja animacji odbicia w~klasie BounceAnimationLogic; \zrodloWlasne},label=lst:animacja]
private async UniTaskVoid PlayBounceAnimation(GameObject gameObject)
{
    CancellationTokenSource cancellationTokenSource = new ();
    UniTask turningRedTask = _animationPlayer.PlayAnimation
        (_ballTurningRedAnimation, _animationDuration, cancellationTokenSource.Token);
    bool collidedObjectIsWall = IsWall(gameObject);
    if (collidedObjectIsWall)
    {
        await _animationPlayer.PlayAnimation
            (_ballSqueezingAnimation, HalfAnimationDuration, cancellationTokenSource.Token);
        await _animationPlayer.PlayAnimation
            (_ballExtendingAnimation, HalfAnimationDuration, cancellationTokenSource.Token);
    }
    await turningRedTask;
    // ... przywrócenie początkowej skali i koloru kuli
    cancellationTokenSource.Dispose();
}
\end{lstlisting}

Ta sama technika posłużyła do zmniejszania zebranych monet i~apteczek, odmierzania czasu
animacji śmierci przeciwnika oraz wypełniania paska ładowania. Kod animacji jest zwykłą
metodą asynchroniczną zapisaną liniowo, bez korutyn i~ręcznego przechowywania stanu.

\section{Kamera i minimapa}

Kamera główna jest sterowana przez wirtualną kamerę pakietu Cinemachine \cite{cinemachine},
której cel śledzenia ustawia komponent \klasa{MainCamera} po otrzymaniu znacznika kuli
z~kontenera. Minimapa korzysta z~drugiej kamery, tworzonej w~kodzie przez klasę
\klasa{MinimapCamera}. Jest to kamera ortograficzna obejmująca obszar 28~$\times$~28~m
(parametr \kod{orthographicSize}, czyli połowa wysokości widoku, wynosi 14~m), umieszczona
40~m nad kulą i~skierowana pionowo w~dół. Renderuje ona obraz do tekstury
256~$\times$~256 pikseli wyświetlanej w~komponencie \klasa{RawImage}. W~każdej klatce kamera minimapy jest
przesuwana nad kulę i~obracana o~ten sam kąt odchylenia co kamera główna, więc góra
minimapy odpowiada kierunkowi, w~którym patrzy gracz.

Kolorowe znaczniki kuli, przeciwników, apteczek i~mety leżą w~osobnej warstwie
\kod{Minimap}, usuwanej z~maski kamery głównej. Komponent \klasa{MinimapIcon} utrzymuje
znacznik poziomo nad właścicielem i~kompensuje jego skalę, a~znacznik przeciwnika znika
razem z~nim, ponieważ jest jego obiektem potomnym.

\section{Pomiar czasu i statystyki rozgrywki}

Czas przejścia poziomu mierzy klasa \klasa{LevelTimerLogic} (listing~\ref{lst:czas}),
której metoda \kod{Tick} jest wywoływana co klatkę przez komponent licznika. Dotarcie do
mety wywołuje metodę \kod{Finish}, która zatrzymuje pomiar, porównuje wynik
z~najlepszym zapisanym czasem, zapisuje statystyki i~publikuje wynik we~właściwości
reaktywnej. Utrata zdrowia jedynie zatrzymuje pomiar. Czas jest formatowany jako
\kod{mm:ss.cc} na podstawie liczby pełnych setnych sekundy, co eliminuje błędne wartości
typu \kod{00:60.00} powstające przy zaokrąglaniu i~uniezależnia zapis od ustawień
regionalnych systemu.

\begin{lstlisting}[caption={Zakończenie pomiaru czasu w~klasie LevelTimerLogic; \zrodloWlasne},label=lst:czas]
public void Finish()
{
    if (!IsRunning) return;
    Stop();
    bool isNewBest = !_levelStatistics.TryGetBestTime(out float bestTime) || Elapsed < bestTime;
    if (isNewBest)
    {
        bestTime = Elapsed;
        _levelStatistics.SaveBestTime(bestTime);
    }
    _levelStatistics.MarkCompleted();
    _levelStatistics.Save();
    Result.Value = new LevelResult(Elapsed, bestTime, isNewBest);
}
\end{lstlisting}

Wyniki są przechowywane w~mechanizmie \klasa{PlayerPrefs} \cite{unity-playerprefs}, który
w~systemie Windows zapisuje dane w~rejestrze. Klasa \klasa{LevelStatistics} buduje klucze z~nazwy poziomu
(tabela~\ref{tab:klucze}), dzięki czemu dodanie nowego poziomu nie wymaga zmian w~kodzie.
Łączny czas gry zapisuje komponent \klasa{LevelStatisticsRecorder} przy opuszczeniu
poziomu oraz przy wstrzymaniu aplikacji, dzięki czemu wynik nie zostaje utracony, nawet
jeśli proces zostanie zakończony bez wywołania metod zamykających.

\begin{table}[!htb]
  \centering
  \caption{Klucze statystyk zapisywanych w~PlayerPrefs; \zrodloWlasne}
  \label{tab:klucze}
  \begin{tabular}{l l l}
    \toprule
    \textbf{Klucz} & \textbf{Typ} & \textbf{Znaczenie} \\
    \midrule
    \kod{Best Time <poziom>} & \kod{float} & najlepszy czas przejścia [s] \\
    \kod{Time Played <poziom>} & \kod{float} & łączny czas spędzony na poziomie [s] \\
    \kod{Max Coins <poziom>} & \kod{int} & najwięcej monet zebranych w~jednym podejściu \\
    \kod{Completed <poziom>} & \kod{int} & 1, jeśli poziom został ukończony \\
    \bottomrule
  \end{tabular}
\end{table}

\begin{figure}[!htb]
  \centering
  \begin{minipage}{0.49\textwidth}
    \centering
    \zrzut[\textwidth]{poziom1_wygrana}\\[2pt]
    {\small a)}
  \end{minipage}\hfill
  \begin{minipage}{0.49\textwidth}
    \centering
    \zrzut[\textwidth]{menu_statystyki}\\[2pt]
    {\small b)}
  \end{minipage}
  \caption{a) Okno końca poziomu po dotarciu do mety z~czasem przejścia, b) okno statystyk
    w~menu głównym przed rozegraniem poziomów; \zrodloWlasne}
  \label{rys:wyniki}
\end{figure}

Okno statystyk w~menu (rysunek~\ref{rys:wyniki}b) odczytuje te same klucze przy każdym
otwarciu. Okno końca poziomu jest wspólne dla wygranej i~przegranej; panel z~czasem
przejścia jest w~nim widoczny tylko wtedy, gdy wynik poziomu został opublikowany, czyli
po dotarciu do mety (rysunek~\ref{rys:wyniki}a).

\section{Przepływ scen, usługi zewnętrzne i walidacja konfiguracji}

Przejścia pomiędzy scenami realizuje komponent \klasa{SceneLoader}. W~edytorze scena
docelowa jest wybierana w~inspektorze jako zasób typu \klasa{SceneAsset}, a~metoda
\kod{OnValidate} zapisuje jej nazwę w~polu serializowanym. W~wersji zbudowanej typ
\klasa{SceneAsset} nie istnieje, więc wykorzystywana jest wyłącznie zapisana nazwa.
Rozwiązuje to częsty problem literówek w~nazwach scen wpisywanych ręcznie. Przycisk
poziomu w~menu zapisuje nazwę poziomu w~obiekcie \klasa{LoadLevel} z~kontekstu projektu
i~wczytuje scenę ładowania, która losuje czas oczekiwania od 1 do 3~s, wypełnia pasek
postępu i~wczytuje zapisany poziom (rysunek~\ref{rys:menu}).

\begin{figure}[!htb]
  \centering
  \begin{minipage}{0.49\textwidth}
    \centering
    \zrzut[\textwidth]{menu_glowne}\\[2pt]
    {\small a)}
  \end{minipage}\hfill
  \begin{minipage}{0.49\textwidth}
    \centering
    \zrzut[\textwidth]{ekran_ladowania}\\[2pt]
    {\small b)}
  \end{minipage}
  \caption{a) Menu główne z~przyciskami poziomów i~odbijającą się piłką w~tle, b) ekran
    ładowania poziomu; \zrodloWlasne}
  \label{rys:menu}
\end{figure}

Scena rozruchowa zawiera komponent \klasa{FirebaseInitializer}, który wywołuje
asynchroniczną metodę \kod{FirebaseApp.CheckAndFixDependenciesAsync}, a~wynik obsługuje
w~wątku głównym; brak wymaganych zależności jest rejestrowany w~dzienniku i~nie przerywa
działania gry \cite{firebase-unity}. Obiekt inicjalizujący nie jest niszczony przy zmianie
sceny. Przygotowano również klasę \klasa{RewardedAd}, która inicjalizuje SDK Appodeal dla
reklam z~nagrodą i~implementuje interfejs zdarzeń reklamy \cite{appodeal}. Klasa ta nie
została jeszcze podłączona do żadnej sceny, a~sposób nagradzania gracza pozostaje
elementem dalszego rozwoju.

Skrypt edytora implementujący interfejs \klasa{IPreprocessBuildWithReport} jest wywoływany
przed każdym budowaniem aplikacji. Sprawdza on, czy wszystkie sceny istnieją, i~waliduje
pliki konfiguracyjne (listing~\ref{lst:walidacja}). Walidacja wyszukuje przez refleksję
wszystkie klasy dziedziczące po \klasa{AbstractConfiguration<T>} i~wywołuje ich metodę
\kod{Load}, więc brak pliku, sekcji lub wartości przerywa budowanie, zamiast ujawnić się
dopiero w~trakcie gry.

\begin{lstlisting}[caption={Walidacja plików konfiguracyjnych przed budowaniem; \zrodloWlasne},label=lst:walidacja]
private static void ValidateConfigurations()
{
    // Instantiate every JSON configuration so missing sections, resources or values fail before building.
    foreach (Type type in typeof(BallMovementConfiguration).Assembly.GetTypes()
                 .Where(t => !t.IsAbstract && t.BaseType != null && t.BaseType.IsGenericType &&
                             t.BaseType.GetGenericTypeDefinition() == typeof(AbstractConfiguration<>)))
    {
        object configuration = Activator.CreateInstance(type);
        try
        {
            type.GetMethod("Load").Invoke(configuration, null);
        }
        catch (Exception exception)
        {
            throw new BuildFailedException($"Invalid configuration {type.Name}: {exception.GetBaseException().Message}");
        }
    }
}
\end{lstlisting}
